\begin{sectionbox}
    \subsection{LTI-System}
    Ein System ist ein LTI-System (Linear, Time Invariant), wenn für beliebige Eingangssignale $x_1(t), x_2(t)$ und Ausgänge $y_1(t), y_2(t)$ gilt:
    \begin{description}
        \item[Linearität] $\mathcal{T}(ax_1(t) + bx_2(t)) = ay_1(t) + by_2(t)$
        \item[Zeitinvarianz] $\mathcal{T}(x(t-t_0))=y(t-t_0)$
    \end{description}

    \subsubsection{Leistung eines periodischen Signals}
    \begin{equation*}
        P=\frac{1}{RT_0}\int_0^{T_0}x^2(t)\diff t=U_{eff}I_{eff}\cos\varphi \quad E=PT_0
    \end{equation*}

    \subsubsection{Effektivwerte (RMS) eines periodischen Signals}
    \begin{equation*}
        X_{\ir eff} = \sqrt{\frac{1}{T_0}\int_0^{T_0}x^2(t)\diff t}
    \end{equation*}

    \subsubsection{Periode eines Signals}
    Größter gemeinsamer Teiler der Frequenz oder kleinstes gemeinsames Vielfaches der Periodendauer
\end{sectionbox}